\subsection{積分変数の選択}

細い帯は縦に取る必要はない.\,$c\leq y\leq d$ において,右側の境界が $x=q(y)$,左側が $x=p(y)$ と表せるなら,帯の横の長さは $q(y)-p(y)$ なので
\begin{equation}\label{式:左右の差の面積}
S=\int_c^d\{q(y)-p(y)\}\dd y
\end{equation}
である.式\eqref{式:上下の差の面積}では上下の差,この式では左右の差を使う.変数の文字を交換するだけでなく,実際にどの向きに帯を取るかが変わっている.

たとえば $x=y^2$ と $x=2-y$ に囲まれる部分では,交点の $y$ 座標は $-2,1$ である.その間では直線が右側なので,
\[
S=\int_{-2}^1(2-y-y^2)\dd y
 =\left[2y-\frac{y^2}{2}-\frac{y^3}{3}\right]_{-2}^1
 =\frac92.
\]
縦に切ると放物線の上側と下側を $y=\pm\sqrt x$ に分け,さらに直線との関係も場合分けすることになる.横に切れば,一つの区間で一つの差になる.境界をどちらの変数で表すと少ない場合分けで済むかが,積分変数を選ぶ基準になる.



